\section{Bảng siêu tham số tối ưu của từng mô hình}
\label{sec:phu-luc-sieu-tham-so}

\begin{table}[H]
\centering
\caption{Siêu tham số tối ưu sau tuning (Mục~\ref{sec:tuning-classic-ml}).}
\label{tab:phu-luc-sieu-tham-so}
\begin{tabular}{@{}lll@{}}
\toprule
\textbf{Mô hình} & \textbf{Siêu tham số} & \textbf{Giá trị tối ưu} \\
\midrule
\TODO{} & \TODO{} & \TODO{} \\
\bottomrule
\end{tabular}
\end{table}

\section{Kết quả chi tiết theo từng fold}
\label{sec:phu-luc-ket-qua-fold}

\TODO{Trình bày bảng kết quả (Macro F1, PR-AUC) của mỗi mô hình trên từng fold
rolling-origin đã mô tả ở Mục~\ref{sec:giao-thuc-chia-du-lieu}, kèm trung bình và độ lệch
chuẩn qua các seed (Mục~\ref{sec:ky-thuat-huan-luyen-dl}).}

\section{Kết quả SAITS theo seed trên tập test cố định}
\label{sec:phu-luc-saits-seed}

Bảng~\ref{tab:saits-theo-seed} trình bày Macro F1 của từng lượt và epoch có
validation weighted cross-entropy thấp nhất. Ba seed dùng cùng tập test
2015--2016 và cùng mask cho một điều kiện; chúng không phải ba fold thời gian.
Trung bình và độ lệch chuẩn mẫu ở Bảng~\ref{tab:saits-macro-f1-missingness}
được tính từ số liệu đầy đủ trước khi làm tròn bốn chữ số thập phân.

\begin{table}[H]
\centering
\caption{Macro F1 của SAITS theo seed. Cột epoch liệt kê checkpoint được chọn
theo thứ tự seed 1/2/42, không phải tổng số epoch đã huấn luyện.}
\label{tab:saits-theo-seed}
\begin{tabular}{@{}lcccc@{}}
\toprule
\textbf{Điều kiện} & \textbf{Seed 1} & \textbf{Seed 2} & \textbf{Seed 42}
& \shortstack{\textbf{Epoch tốt nhất}\\\textbf{(1/2/42)}} \\
\midrule
% Generated from report/data/saits_verified_results.json; SAITS only.
Không ẩn thêm & $0.9130$ & $0.9278$ & $0.9264$ & 12/22/27 \\
Point 10\% & $0.9310$ & $0.9243$ & $0.9246$ & 25/19/37 \\
Point 30\% & $0.8277$ & $0.8023$ & $0.8301$ & 13/25/10 \\
Point 50\% & $0.8510$ & $0.8609$ & $0.8471$ & 13/19/10 \\
Subsequence 10\% & $0.9115$ & $0.9011$ & $0.9078$ & 17/17/31 \\
Subsequence 30\% & $0.8750$ & $0.8629$ & $0.8665$ & 49/33/19 \\
Subsequence 50\% & $0.8411$ & $0.8357$ & $0.7321$ & 16/36/8 \\
Block 10\% & $0.9088$ & $0.9235$ & $0.9200$ & 11/18/15 \\
Block 30\% & $0.9107$ & $0.8992$ & $0.9156$ & 25/18/38 \\
Block 50\% & $0.8909$ & $0.9014$ & $0.8906$ & 15/17/38 \\
\bottomrule

\end{tabular}
\end{table}

\paragraph{Nguồn số liệu và kiểm tra nhất quán.}
Thí nghiệm gốc có mã \texttt{experiment\_a2a9e03b4630}, được thực hiện bằng
notebook \path{notebooks/SAITS_Google_Colab.ipynb}. Đợt đối soát ngày
24/11/2026 đã kiểm tra 30 lượt hoàn tất, tính lại Macro F1 từ 60 file dự đoán
validation/test, đối chiếu với metrics từng lượt và bảng tổng hợp, đồng thời
kiểm tra 180 checksum của các artifact. Các giá trị Macro F1 tính lại khớp
với số liệu đã lưu. Bản số liệu SAITS dùng để tạo báo cáo được lưu tại
\path{report/data/saits_verified_results.json}, gồm cấu hình, kết quả từng
seed và số liệu tổng hợp; bản này được trích từ
\path{outputs/saits_comparison_20261124/audit.json}.

Bảng và hình riêng của SAITS có thể được tái tạo bằng
\texttt{python report/make\_saits\_results\_figure.py}. Script đọc bản số liệu
đã kiểm chứng, kiểm tra đủ ba seed cho mỗi điều kiện, tính sample std với
mẫu số $3-1$ và xuất các hàng bảng cùng hình
\path{report/figs/saits_macro_f1_missingness.png}. Cấu hình ở
Bảng~\ref{tab:saits-cau-hinh} là cấu hình thực nghiệm cố định, chưa được xác
nhận là cấu hình tối ưu qua tìm kiếm siêu tham số.

\section{Phân công thành viên}
\label{sec:phu-luc-phan-cong}

\begin{table}[H]
\centering
\caption{Phân công công việc giữa các thành viên.}
\label{tab:phan-cong-thanh-vien}
\begin{tabular}{@{}lll@{}}
\toprule
\textbf{Thành viên} & \textbf{Công việc phụ trách} & \textbf{Tỷ lệ đóng góp} \\
\midrule
\TODO{Họ và tên} & \TODO{} & \TODO{} \\
\TODO{Họ và tên} & \TODO{} & \TODO{} \\
\bottomrule
\end{tabular}
\end{table}

\section{Link mã nguồn}
\label{sec:phu-luc-link-ma-nguon}

\TODO{Bổ sung đường dẫn repository mã nguồn (ví dụ GitHub) và hướng dẫn chạy lại pipeline
(tham khảo README.md ở gốc repository).}
